\section{问题重述}

霍尔木兹海峡连接波斯湾和阿曼湾，承担全球原油海运的重要份额，是中东原油进入国际市场的关键通道。若该通道因地缘冲突、航运安全或军事风险发生封锁，国际原油市场将同时面对实际供应缺口、运输成本上升、市场恐慌放大和政策缓冲响应，价格可能在短期内剧烈波动，并在封锁延长后进入新的风险定价区间。

赛题要求基于布伦特原油期货价格数据和国际原油供需规律，建立数学模型分析霍尔木兹海峡封锁对国际原油价格的影响。结合题意，本文需要回答四个问题：

\begin{enumerate}
  \item 构建短期冲击模型，量化 0--60 天内封锁对油价的即时影响，解释供应缺口、政策缓冲和风险溢价的共同作用。
  \item 构建中长期油价调节模型，预测封锁持续 60--180 天时的价格演化路径和可能的价格平衡区间。
  \item 识别影响油价的关键参数，分析供应中断、SPR 释放、商业库存、绕道运输、需求弹性和地缘风险变量的敏感性。
  \item 检验模型稳定性，给出可用于风险预警和政策评估的量化结论。
\end{enumerate}

\begin{table}[!htbp]
\centering
\caption{赛题任务与本文模型回应}
\begin{tabularx}{\linewidth}{lX X}
\toprule
赛题要求 & 题面重点 & 本文回应 \\
\midrule
任务 1：短期冲击模型（0--60 天） & 建立 0--60 天日度供需平衡模型，使用附件布伦特原油价格数据，并考虑供应中断、SPR、库存、绕道、恐慌需求和短期弹性。 & 完成附件数据清洗、传统供需基准、短期冲击模型、参数校准和鲁棒性检验，给出 RMSE、MAE、MAPE、基准对比、滞后检验、机制消融和参数扰动测试。 \\
任务 2：中长期油价调节模型（60--180 天） & 若封锁持续到 60--180 天，预测油价变化和价格平衡点，考虑增产、长期弹性、阿联酋增产约束和库存耗尽跳变风险。 & 构建乐观、中性、悲观三情景路径，给出第 60/90/120/180 天价格、外推期峰值、库存风险、二次跳涨风险和关键参数敏感性排序。 \\
\bottomrule
\end{tabularx}
\end{table}
